\documentclass[10pt,a4paper]{amsart}
\usepackage[margin=19mm]{geometry}
\usepackage{amsmath,amssymb,mathtools}
\usepackage[scr=rsfs,cal=euler]{mathalfa}
\usepackage{xfrac}
\usepackage[unicode,hidelinks,bookmarks=false]{hyperref}
\newcommand{\ie}{i.e.\ }
\newcommand{\cf}{cf.\ }
\newcommand{\resp}{respectively\ }
\newcommand{\Subsection}{\subsection}
\providecommand{\nobreakdash}{\nobreak\hbox{-}}
\setlength{\emergencystretch}{2em}
\multlinegap0pt
\usepackage{mathtools}
\usepackage{multirow}
\usepackage{bm}
\usepackage{amssymb}
\usepackage{bbm}
\usepackage{dsfont}
\usepackage[all]{xy}
\usepackage{tikz}
\newtheorem{corollary}{Corollary}
\newtheorem{lemma}{Lemma}
\newcommand{\Fac}{\opname{Fac}}
\def\mod{\opname{mod}\nolimits}
\newcommand{\opname}[1]{\operatorname{\mathsf{#1}}}
\title{Newton polytopes in cluster algebras and $\tau$-tilting theory}
\author{Peigen Cao}
\date{Source: arXiv:2602.07929v2 (CC BY 4.0)}
\begin{document}
\maketitle

\noindent\emph{Source and licence.} Newton polytopes in cluster algebras and $\tau$-tilting theory. Version arXiv:2602.07929v2 (CC BY 4.0), distributed by arXiv under the Creative Commons Attribution 4.0 International licence, http://creativecommons.org/licenses/by/4.0/, as recorded in the arXiv licence metadata for this exact version and shown on the abstract page https://arxiv.org/abs/2602.07929v2. Full licence notice: \texttt{licenses/arxiv-2602.07929-CC-BY-4.0.txt}.\par\smallskip

\noindent\textit{Editorial scope.} Counted unit: the article's lemma lem:UV-torsion (source0.tex lines 1509--1511). The article's complete original proof of it is delivered with it (source0.tex lines 1512--1520, one \texttt{proof} environment of 9 lines). No mathematics is written, repaired or re-derived: the statement and the proof are the article's own lines, in the article's own notation, and no retained line is substituted or re-flowed.  Retained uncounted as the source-local context the proof needs: lines 1339--1341.  Nothing was omitted from any retained span, so the package carries no omission record.  No retained line contains a citation, so the wrapper carries no bibliography and no bibliography entry.  No retained line contains a figure environment, an image-inclusion command or drawing code, so no image is delivered and the package declares no asset dependency.  Editorial formatting only, in addition: an AMS \texttt{amsart} preamble that restates the article's own declarations and macros used by the retained lines, namely the environments \texttt{corollary}, \texttt{lemma} on separate counters, together with the standard packages those lines need.  The retained environments print in this document as Corollary 1, Lemma 1 in order of appearance, whereas the article prints the same items with the numbering of its own full document.  The complete native source of the article as distributed in the arXiv e-print for this exact version is bundled unchanged as \texttt{sources/194220/source0.tex}.
\begin{corollary}\label{cor:dom-hom}
Let $\mathcal M=(M,P)$ be a $\tau$-rigid pair in $\mod A$. Then a module $N\in \mod A$ belongs to   $\prescript{\bot}{}(\tau M)\cap P^{\bot}$ if and only if $F_N[{\bf g}_{\mathcal M}]=0$, that is, $\langle{\bf g}_{\mathcal M},\underline{\dim} N_0\rangle\leq 0$ for any quotient module $N_0$ of $N$.
\end{corollary}

\begin{lemma}\label{lem:UV-torsion}
    Let $U$ and $V$ be two left finite modules in $\mod A$. If $U$ and $V$ have the same Newton polytope, then $\langle U\rangle_{\rm tors}=\langle V\rangle_{\rm tors}$.
\end{lemma}

\begin{proof}
 Since $U$ and $V$ are left finite modules, there are two basic $\tau$-tilting pairs $\mathcal M=(M,P)$ and $\mathcal N=(N, Q)$ such that $\Fac M=\langle U\rangle_{\rm tors}$ and $\Fac N=\langle V\rangle_{\rm tors}$.
 Since $U\in\Fac M=\prescript{\bot}{}(\tau M)\cap P^{\bot}$, $V\in\Fac N=\prescript{\bot}{}(\tau N)\cap Q^{\bot}$ and
 by Corollary \ref{cor:dom-hom}, we have $F_U[{\bf g}_{\mathcal M}]=0$ and $F_V[{\bf g}_{\mathcal N}]=0$.

 Since $U$ and $V$ have the same Newton polytope, we know that $F_U[{\bf r}]=F_V[{\bf r}]$ for any ${\bf r}\in\mathbb Z^n$. In particular, we have $F_U[{\bf g}_{\mathcal N}]=F_V[{\bf g}_{\mathcal N}]=0$ and $F_V[{\bf g}_{\mathcal M}]=F_U[{\bf g}_{\mathcal M}]=0$.

 Then by Corollary \ref{cor:dom-hom} again, we see that $U\in \prescript{\bot}{}(\tau N)\cap Q^{\bot}=\Fac N=\langle V\rangle_{\rm tors}$ and $V\in \prescript{\bot}{}(\tau M)\cap P^{\bot}=\Fac M=\langle U\rangle_{\rm tors}$. Hence, we have $\langle U\rangle_{\rm tors}=\langle V\rangle_{\rm tors}$.
\end{proof}

\end{document}
